\section{Background and Motivation}
\textbf{Speculative Decoding.}
Large language models (LLMs) generate text auto-regressively: given a prompt $(x_1, \ldots, x_n)$, they produce an output sequence $(x_{n+1}, \ldots, x_{n+S})$, 
one token at a time. At each step, the model computes a probability distribution over the vocabulary and samples the next token. 
% This strictly sequential decoding limits the utilization of highly parallel hardware such as GPUs~\cite{pope2023efficiently} and results in high latency. 
However, while LLMs \emph{generate} tokens sequentially, they can also \emph{evaluate} the probabilities of a \emph{given} token sequence in parallel. That is, for candidate tokens $x_{n+1}, \ldots, x_{n+S}$, an LLM can compute $P(x_{n+1} \mid x_1, \ldots, x_n), \ldots, P(x_{n+S} \mid x_1, \ldots, x_{n+S-1})$ in a single forward pass.
Speculative decoding~\cite{leviathan2023fast, chen2023accelerating} utilizes this property of LLMs as an evaluator on a string of candidate tokens in parallel, enabling higher hardware utilization and potentially generating multiple tokens.
For example, a smaller model may generate $k$ candidate tokens, and our target LLM model then predicts $k+1$ probability distributions for each token position in parallel. We then accept the first $m$ tokens from the $k$ candidate tokens that satisfy the sampling method. 

\textbf{SD Variations.}
When SD was first introduced, the standard approach was to use a smaller draft model with the same vocabulary as the target model, since rejection sampling must operate on a probability distribution in the same token embedding space.
Subsequent work has proposed various ways to construct the draft model, such as quantization~\cite{lin2023awq, frantar2023gptq, kim2024squeezellm} or distillation~\cite{zhou2023distillspec} of the target LLM.
However, in practice, an off-the-shelf draft model is not always available, requiring quantizing or pruning the original model, or even pretraining a smaller model from scratch.

To overcome the limitations of requiring a separate draft model, \emph{draft-model-free} SD methods have been proposed.
One common approach is to add auxiliary prediction layers on top of the final transformer block ~\cite{li2024eagle, li2024eagle2fasterinferencelanguage, li2025eagle3}.
These additional ``head'' layers take the hidden states as input and directly propose tokens.
However, practitioners still need to fine-tune these drafting heads separately.
% Another line of work fully eliminates the reuqirement of training. \emph{Multi-Token Prediction (MTP)}~\cite{deepseekai2025deepseekv3technicalreport, 5team2025glm45agenticreasoningcoding} co-trains the auxiliary heads jointly with the main model, rather than fine-tuning them afterwards. As a result, auxiliary heads are provided out-of-the-box with the model, while also achieving a significantly higher token acceptance rate compared to the post-training, add-on approaches.
Recently, a new line of work has removed the need for post-hoc training. Multi-Token Prediction (MTP)~\cite{deepseekai2025deepseekv3technicalreport,5team2025glm45agenticreasoningcoding} co-trains auxiliary prediction heads jointly with the main model, rather than fine-tuning them afterward. As a result, these auxiliary heads are available out of the box and achieve substantially higher token-acceptance rates.
% highlighting the potential of training-time co-design \alvin{what's training time co design?} for speculative decoding.

Another line of draft-model-free SD methods utilize a non-LLM proposer.
A representative example is prompt lookup decoding, or \emph{$n$-gram}~\cite{saxena2023prompt}.
The intuition is that LLM outputs often contain recurring phrases; therefore, we can extract $n$-gram snippets from previously generated text and reuse them as token proposals.
This approach is particularly effective in workloads such as code or passage editing, where the input and output share substantial overlap. 
Another example is using retrieval-augmented generation (RAG) to retrieve candidate proposals from an external corpus~\cite{he2024restretrievalbasedspeculativedecoding}.
These non-LLM approaches demonstrate the potential for lightweight or task-specific drafting mechanisms.
% However, it remains an open problem to determine how to construct the most effective non-LLM proposers, and how to combine them optimally with standard SD pipelines.

\textbf{Problems of Existing Benchmarks.}
Most prior works evaluate SD by building proof-of-concept research prototypes, as integration into a production-grade inference system is nontrivial and requires modifications across multiple components.
However, this makes a fair comparison between SD methods difficult and the benefits in real-world deployment impossible to predict, due to the following limitations.
First, many prototype implementations lack key optimizations in production systems, such as the use of CUDA graphs~\cite{cudagraph} to reduce system overhead, a common feature in LLM serving systems~\cite{kwon2023efficient, zheng2024sglang}.
Second, the implementations are typically evaluated with a batch size of one, which is unrepresentative of large-scale batched deployments of the real world.
Finally, prior work lacks an in-depth analysis of SD speedup, as they primarily focus on high-level metrics such as average latency or dataset-level acceptance rate. \cite{li2024eagle,li2025eagle3,xia-etal-2024-unlocking}
We believe that a deeper understanding, such as execution time breakdown and position-level acceptance behavior, is necessary to fully explain and optimize the end-to-end performance of SD.